% Options for packages loaded elsewhere
% Options for packages loaded elsewhere
\PassOptionsToPackage{unicode}{hyperref}
\PassOptionsToPackage{hyphens}{url}
\PassOptionsToPackage{dvipsnames,svgnames,x11names}{xcolor}
%
\documentclass[
  spanish,
  letterpaper,
  DIV=11,
  numbers=noendperiod]{scrartcl}
\usepackage{xcolor}
\usepackage{amsmath,amssymb}
\setcounter{secnumdepth}{5}
\usepackage{iftex}
\ifPDFTeX
  \usepackage[T1]{fontenc}
  \usepackage[utf8]{inputenc}
  \usepackage{textcomp} % provide euro and other symbols
\else % if luatex or xetex
  \usepackage{unicode-math} % this also loads fontspec
  \defaultfontfeatures{Scale=MatchLowercase}
  \defaultfontfeatures[\rmfamily]{Ligatures=TeX,Scale=1}
\fi
\usepackage{lmodern}
\ifPDFTeX\else
  % xetex/luatex font selection
\fi
% Use upquote if available, for straight quotes in verbatim environments
\IfFileExists{upquote.sty}{\usepackage{upquote}}{}
\IfFileExists{microtype.sty}{% use microtype if available
  \usepackage[]{microtype}
  \UseMicrotypeSet[protrusion]{basicmath} % disable protrusion for tt fonts
}{}
\makeatletter
\@ifundefined{KOMAClassName}{% if non-KOMA class
  \IfFileExists{parskip.sty}{%
    \usepackage{parskip}
  }{% else
    \setlength{\parindent}{0pt}
    \setlength{\parskip}{6pt plus 2pt minus 1pt}}
}{% if KOMA class
  \KOMAoptions{parskip=half}}
\makeatother
% Make \paragraph and \subparagraph free-standing
\makeatletter
\ifx\paragraph\undefined\else
  \let\oldparagraph\paragraph
  \renewcommand{\paragraph}{
    \@ifstar
      \xxxParagraphStar
      \xxxParagraphNoStar
  }
  \newcommand{\xxxParagraphStar}[1]{\oldparagraph*{#1}\mbox{}}
  \newcommand{\xxxParagraphNoStar}[1]{\oldparagraph{#1}\mbox{}}
\fi
\ifx\subparagraph\undefined\else
  \let\oldsubparagraph\subparagraph
  \renewcommand{\subparagraph}{
    \@ifstar
      \xxxSubParagraphStar
      \xxxSubParagraphNoStar
  }
  \newcommand{\xxxSubParagraphStar}[1]{\oldsubparagraph*{#1}\mbox{}}
  \newcommand{\xxxSubParagraphNoStar}[1]{\oldsubparagraph{#1}\mbox{}}
\fi
\makeatother


\usepackage{longtable,booktabs,array}
\newcounter{none} % for unnumbered tables
\usepackage{calc} % for calculating minipage widths
% Correct order of tables after \paragraph or \subparagraph
\usepackage{etoolbox}
\makeatletter
\patchcmd\longtable{\par}{\if@noskipsec\mbox{}\fi\par}{}{}
\makeatother
% Allow footnotes in longtable head/foot
\IfFileExists{footnotehyper.sty}{\usepackage{footnotehyper}}{\usepackage{footnote}}
\makesavenoteenv{longtable}
\usepackage{graphicx}
\makeatletter
\newsavebox\pandoc@box
\newcommand*\pandocbounded[1]{% scales image to fit in text height/width
  \sbox\pandoc@box{#1}%
  \Gscale@div\@tempa{\textheight}{\dimexpr\ht\pandoc@box+\dp\pandoc@box\relax}%
  \Gscale@div\@tempb{\linewidth}{\wd\pandoc@box}%
  \ifdim\@tempb\p@<\@tempa\p@\let\@tempa\@tempb\fi% select the smaller of both
  \ifdim\@tempa\p@<\p@\scalebox{\@tempa}{\usebox\pandoc@box}%
  \else\usebox{\pandoc@box}%
  \fi%
}
% Set default figure placement to htbp
\def\fps@figure{htbp}
\makeatother



\ifLuaTeX
\usepackage[bidi=basic,shorthands=off]{babel}
\else
\usepackage[bidi=default,shorthands=off]{babel}
\fi
\ifLuaTeX
  \usepackage{selnolig} % disable illegal ligatures
\fi


\setlength{\emergencystretch}{3em} % prevent overfull lines

\providecommand{\tightlist}{%
  \setlength{\itemsep}{0pt}\setlength{\parskip}{0pt}}



 


\KOMAoption{captions}{tableheading}
\makeatletter
\@ifpackageloaded{caption}{}{\usepackage{caption}}
\AtBeginDocument{%
\ifdefined\contentsname
  \renewcommand*\contentsname{Tabla de contenidos}
\else
  \newcommand\contentsname{Tabla de contenidos}
\fi
\ifdefined\listfigurename
  \renewcommand*\listfigurename{Listado de Figuras}
\else
  \newcommand\listfigurename{Listado de Figuras}
\fi
\ifdefined\listtablename
  \renewcommand*\listtablename{Listado de Tablas}
\else
  \newcommand\listtablename{Listado de Tablas}
\fi
\ifdefined\figurename
  \renewcommand*\figurename{Figura}
\else
  \newcommand\figurename{Figura}
\fi
\ifdefined\tablename
  \renewcommand*\tablename{Tabla}
\else
  \newcommand\tablename{Tabla}
\fi
}
\@ifpackageloaded{float}{}{\usepackage{float}}
\floatstyle{ruled}
\@ifundefined{c@chapter}{\newfloat{codelisting}{h}{lop}}{\newfloat{codelisting}{h}{lop}[chapter]}
\floatname{codelisting}{Listado}
\newcommand*\listoflistings{\listof{codelisting}{Listado de Listados}}
\makeatother
\makeatletter
\makeatother
\makeatletter
\@ifpackageloaded{caption}{}{\usepackage{caption}}
\@ifpackageloaded{subcaption}{}{\usepackage{subcaption}}
\makeatother
\usepackage{bookmark}
\IfFileExists{xurl.sty}{\usepackage{xurl}}{} % add URL line breaks if available
\urlstyle{same}
% fallback for those not using the hyperref driver hyperxmp:
\makeatletter
\@ifundefined{xmpquote}{\newcommand{\xmpquote}[1]{#1}}{}
\makeatother
\hypersetup{
  pdftitle={Análisis del contexto nacional y global},
  pdfauthor={Equipo: por definir},
  pdflang={es},
  colorlinks=true,
  linkcolor={blue},
  filecolor={Maroon},
  citecolor={Blue},
  urlcolor={Blue},
  pdfcreator={LaTeX via pandoc}}


\title{Análisis del contexto nacional y global}
\usepackage{etoolbox}
\makeatletter
\providecommand{\subtitle}[1]{% add subtitle to \maketitle
  \apptocmd{\@title}{\par {\large #1 \par}}{}{}
}
\makeatother
\subtitle{Transformación digital y empleo}
\author{Equipo: por definir}
\date{2026-07-31}
\begin{document}
\maketitle

\renewcommand*\contentsname{Tabla de contenidos}
{
\setcounter{tocdepth}{3}
\tableofcontents
}

\section{Resumen ejecutivo}\label{resumen-ejecutivo}

Este informe analiza descriptivamente la relación entre conectividad
digital y empleo en Ecuador, Argentina, Colombia y Venezuela durante
2020--2024. Se emplean 80 observaciones país-año de indicadores del
Banco Mundial (WDI); el indicador de uso de internet procede
originalmente de la Unión Internacional de Telecomunicaciones (UIT). La
validación no detectó duplicados ni valores faltantes en el periodo
analítico.

Entre 2020 y 2024 aumentaron el uso de internet y la tasa de
empleo/población en los cuatro países. Una regresión exploratoria con
efectos fijos por país estima una asociación positiva entre uso de
internet y empleo, pero no identifica un efecto causal. La muestra
contiene solo 20 observaciones para el modelo y cuatro países; por ello
los resultados deben leerse como evidencia descriptiva y exploratoria.

\section{Planteamiento del problema y
objetivos}\label{planteamiento-del-problema-y-objetivos}

La pregunta de investigación es: ¿cómo se relaciona el avance de la
conectividad digital con la tasa de empleo de la población de 15 años o
más en Ecuador, Argentina, Venezuela y Colombia durante 2020--2024?

El objetivo general es analizar esa relación y comparar la evolución de
Ecuador con los tres países seleccionados. Los objetivos específicos son
describir los indicadores digitales y laborales, realizar la comparación
regional y estimar una asociación exploratoria controlando por el
crecimiento real del PIB per cápita.

\section{Fuentes y metodología}\label{fuentes-y-metodologuxeda}

La fuente de descarga es
\href{https://databank.worldbank.org/source/world-development-indicators}{World
Development Indicators del Banco Mundial}. Se utilizaron: tasa de
empleo/población, 15+ (\texttt{SL.EMP.TOTL.SP.ZS}); crecimiento anual
del PIB real per cápita (\texttt{NY.GDP.PCAP.KD.ZG}); desempleo total
(\texttt{SL.UEM.TOTL.ZS}); y uso de internet (\texttt{IT.NET.USER.ZS}).
Este último tiene como fuente primaria la
\href{https://datahub.itu.int/}{UIT DataHub}.

Las respuestas originales se conservaron en \texttt{data/raw/} con un
manifiesto de URL y huellas SHA-256. El periodo 2025 se excluyó del
análisis porque el indicador de uso de internet no tenía valores para
los cuatro países. No se imputaron ni corrigieron valores.

La variable de resultado es la tasa de empleo/población de 15 años o
más. La variable explicativa principal es el porcentaje de personas que
usan internet; el crecimiento del PIB real per cápita es un control. El
desempleo se usa solo como indicador descriptivo por su relación
contable con el empleo.

\section{Procesamiento y
validación}\label{procesamiento-y-validaciuxf3n}

El conjunto procesado contiene 80 observaciones: cuatro países, cinco
años y cuatro indicadores. La validación de datos brutos verificó
integridad por SHA-256, esquema, claves país-año, duplicados, cobertura
y faltantes. Para 2020--2024 no hubo valores nulos ni claves duplicadas.
El reporte completo está en \texttt{outputs/reports/data\_quality.md}.

\section{Resultados descriptivos}\label{resultados-descriptivos}

{\def\LTcaptype{none} % do not increment counter
\begin{longtable}[]{@{}
  >{\raggedright\arraybackslash}p{(\linewidth - 12\tabcolsep) * \real{0.1111}}
  >{\raggedleft\arraybackslash}p{(\linewidth - 12\tabcolsep) * \real{0.1481}}
  >{\raggedleft\arraybackslash}p{(\linewidth - 12\tabcolsep) * \real{0.1481}}
  >{\raggedleft\arraybackslash}p{(\linewidth - 12\tabcolsep) * \real{0.1481}}
  >{\raggedleft\arraybackslash}p{(\linewidth - 12\tabcolsep) * \real{0.1481}}
  >{\raggedleft\arraybackslash}p{(\linewidth - 12\tabcolsep) * \real{0.1481}}
  >{\raggedleft\arraybackslash}p{(\linewidth - 12\tabcolsep) * \real{0.1481}}@{}}
\toprule\noalign{}
\begin{minipage}[b]{\linewidth}\raggedright
País
\end{minipage} & \begin{minipage}[b]{\linewidth}\raggedleft
Internet 2020
\end{minipage} & \begin{minipage}[b]{\linewidth}\raggedleft
Internet 2024
\end{minipage} & \begin{minipage}[b]{\linewidth}\raggedleft
Cambio (pp)
\end{minipage} & \begin{minipage}[b]{\linewidth}\raggedleft
Empleo 2020
\end{minipage} & \begin{minipage}[b]{\linewidth}\raggedleft
Empleo 2024
\end{minipage} & \begin{minipage}[b]{\linewidth}\raggedleft
Cambio (pp)
\end{minipage} \\
\midrule\noalign{}
\endhead
\bottomrule\noalign{}
\endlastfoot
Argentina & 85.51 & 89.67 & 4.15 & 49.60 & 57.44 & 7.84 \\
Colombia & 69.80 & 79.35 & 9.55 & 53.28 & 57.42 & 4.14 \\
Ecuador & 70.70 & 77.17 & 6.47 & 56.36 & 62.31 & 5.95 \\
Venezuela & 72.44 & 76.68 & 4.24 & 44.55 & 49.28 & 4.73 \\
\end{longtable}
}

Los cuatro países registran aumentos simultáneos de conectividad y
empleo en el periodo. Esta coincidencia no permite atribuir el cambio
laboral únicamente a la digitalización: también intervienen ciclo
económico, políticas públicas, estructura sectorial y condiciones
institucionales.

\section{Análisis econométrico
exploratorio}\label{anuxe1lisis-economuxe9trico-exploratorio}

Se estimó la ecuación:

\[Empleo_{it}=\alpha_i+\beta Internet_{it}+\gamma CrecimientoPIBpc_{it}+\varepsilon_{it}.\]

El modelo usa efectos fijos por país y errores estándar HC3. El
coeficiente estimado para uso de internet fue 0.427 puntos porcentuales
de empleo por cada punto porcentual adicional de uso de internet; el
intervalo aproximado al 95\% fue de 0.085 a 0.770. El coeficiente del
crecimiento del PIB real per cápita fue 0.132. Los resultados completos
están en \texttt{outputs/tables/model\_results.csv}.

No se incluyen efectos fijos de año por el tamaño reducido de la
muestra. Los valores p son orientativos y no sustentan una
interpretación causal.

\section{Riesgos, oportunidades y recomendaciones para
Ecuador}\label{riesgos-oportunidades-y-recomendaciones-para-ecuador}

El aumento observado de uso de internet coincide con una mejora de la
tasa de empleo ecuatoriana en el periodo. Esto es compatible con
mecanismos como mejor acceso a información, mercados y oportunidades de
trabajo, pero no demuestra que uno cause el otro.

Se recomienda ampliar el estudio con más países o información
subnacional y microdatos, incorporar indicadores de calidad y
asequibilidad de conectividad, y distinguir empleo formal, informal,
juvenil y por sexo. Cualquier decisión de política debe complementarse
con evidencia causal y contexto institucional.

\section{Limitaciones}\label{limitaciones}

\begin{itemize}
\tightlist
\item
  Muestra pequeña: cuatro países y cinco años.
\item
  El indicador de conectividad mide uso de internet, no calidad,
  habilidades digitales ni tipo de empleo.
\item
  Las estimaciones laborales modeladas mejoran comparabilidad, pero no
  sustituyen el análisis de fuentes nacionales.
\item
  El estudio documenta asociaciones; no identifica causalidad.
\end{itemize}

\section{Referencias}\label{referencias}

\begin{itemize}
\tightlist
\item
  Banco Mundial. \emph{World Development Indicators}.
  https://databank.worldbank.org/source/world-development-indicators
\item
  Unión Internacional de Telecomunicaciones. \emph{ITU DataHub}.
  https://datahub.itu.int/
\item
  Organización Internacional del Trabajo. \emph{ILOSTAT}.
  https://ilostat.ilo.org/data/bulk/
\end{itemize}

\section{Anexos}\label{anexos}

Los archivos \texttt{data/raw/wdi\_manifest\_2020\_2025.json},
\texttt{data/processed/processing\_manifest.json},
\texttt{outputs/reports/data\_quality.md} y
\texttt{outputs/tables/model\_results.csv} contienen la trazabilidad
técnica del análisis.




\end{document}
